\documentclass{article}
\usepackage[utf8]{inputenc}
\usepackage{graphicx}
\usepackage{geometry}
\usepackage{hyperref}
\geometry{a4paper, margin=1in}
\hypersetup{
    colorlinks=true,
    linkcolor=blue,
    filecolor=magenta,
    urlcolor=cyan,
}

\title{Sistema Hospitalario Basado en Microservicios - Informe Sprint 3}
\author{Equipo \#21 \\ Santos Alejandro Arellano Olarte // A01643742 \\ Roberto Anwar Teigeiro Aguilar // A01643651 \\ Fernando Antonio Lopez Garcia // A01643685 \\ Diego Michell Villa Duran // A00836723 \\ Sadrac Aramburo Enciso // A01643639}
\date{16 de Mayo de 2025}

\begin{document}

\maketitle

\section{Introducción}
Este informe documenta la implementación de un sistema hospitalario basado en microservicios para el Sprint 3, que incluye la gestión de pacientes, agendamiento de citas y alertas médicas, asegurado con JWT (JSON Web Tokens) y utilizando Kafka para la comunicación asincrónica. El sistema está diseñado para ser escalable, seguro y eficiente, cumpliendo con los requisitos establecidos en la actividad.

\section{Arquitectura del Sistema}
El sistema está desplegado en un entorno de Docker, utilizando un enfoque de microservicios. La arquitectura se representa en el siguiente diagrama de despliegue UML:

\begin{figure}[h]
    \centering
    \includegraphics[width=\textwidth]{architecture.png}
    \caption{Diagrama de Despliegue UML}
\end{figure}

\section{Decisiones Técnicas}
A continuación, se detallan las decisiones técnicas tomadas durante la implementación del sistema:

\begin{itemize}
    \item \textbf{Microservicios}: Se implementaron tres microservicios: \textit{Pacientes} para la gestión de pacientes (interno y seguro), \textit{Citas} para el agendamiento de citas (uso público limitado) y \textit{Alertas} para procesar alertas médicas basadas en eventos (reactivo). Cada microservicio utiliza Spring Boot para crear APIs REST que manejan datos en formato JSON.
    \item \textbf{Seguridad}: Se utilizó Keycloak como proveedor de identidad para la autenticación basada en JWT. Cada microservicio valida los tokens JWT proporcionados por Keycloak, asegurando que solo usuarios autorizados accedan a los endpoints. Se definieron roles diferenciados (admin, paciente, médico) para controlar el acceso.
    \item \textbf{Comunicación}: Se emplearon APIs REST para la comunicación sincrónica entre el API Gateway y los microservicios. Para la comunicación asincrónica, se utilizó Kafka para manejar eventos como actualizaciones de pacientes que disparan alertas en el microservicio Alertas. El API Gateway facilita el acceso externo y rutea las solicitudes a los microservicios correspondientes.
    \item \textbf{Base de Datos}: Se utilizó PostgreSQL para almacenar los datos de pacientes y citas. Cada microservicio tiene su propio esquema en la base de datos para mantener la separación de preocupaciones.
    \item \textbf{Despliegue}: El sistema se despliega utilizando Docker Compose, lo que permite un manejo sencillo de contenedores y servicios. Se utilizan volúmenes para persistir datos críticos como la base de datos y los metadatos de Kafka.
\end{itemize}

\section{Resultados de Pruebas}
Se realizaron pruebas exhaustivas para verificar la funcionalidad del sistema, especialmente en lo que respecta a la seguridad con JWT y la comunicación entre componentes. A continuación, se presentan los resultados:

\begin{enumerate}
    \item \textbf{Verificación de Salud de los Servicios}: Se verificaron los endpoints de salud de cada microservicio y el API Gateway para confirmar que estaban operativos.
        \begin{itemize}
            \item \texttt{curl http://localhost:8081/actuator/health} - Respuesta: \texttt{\{"status":"UP"\}}
            \item \texttt{curl http://localhost:8082/actuator/health} - Respuesta: \texttt{\{"status":"UP"\}}
            \item \texttt{curl http://localhost:8083/actuator/health} - Respuesta: \texttt{\{"status":"UP"\}}
            \item \texttt{curl http://localhost:8084/actuator/health} - Respuesta: \texttt{\{"status":"UP"\}}
        \end{itemize}

    \item \textbf{Autenticación con JWT}: Se obtuvo un token JWT utilizando las credenciales de administrador y se probaron accesos a endpoints protegidos.
        \begin{itemize}
            \item Obtener token:
                \begin{verbatim}
                curl -X POST http://localhost:8080/realms/hospital/protocol/openid-connect/token \
                -d "client_id=hospital-client" \
                -d "client_secret=hospital-client-secret" \
                -d "username=admin" \
                -d "password=admin" \
                -d "grant_type=password"
                \end{verbatim}
                Respuesta: Contiene \texttt{access_token}.

            \item Acceso a endpoint protegido con token:
                \begin{verbatim}
                curl -H "Authorization: Bearer <access_token>" http://localhost:8084/api/pacientes
                \end{verbatim}
                Respuesta: Lista JSON de pacientes, e.g., \texttt{[{"id":1,"nombre":"Test Patient"}]}

            \item Acceso a endpoint protegido sin token:
                \begin{verbatim}
                curl http://localhost:8084/api/pacientes
                \end{verbatim}
                Respuesta: \texttt{401 Unauthorized}.
        \end{itemize}

    \item \textbf{Comunicación con Kafka}: Se envió un mensaje a través de Kafka para simular una actualización de paciente y se verificó que el microservicio Alertas lo procesara.
        \begin{itemize}
            \item Producir mensaje:
                \begin{verbatim}
                docker exec -it hospital-microservices-kafka-1 kafka-console-producer \
                --broker-list kafka:29092 --topic pacientes-topic
                # Mensaje: {"id":1,"nombre":"Test Patient","status":"updated"}
                \end{verbatim}

            \item Verificar logs de Alertas:
                \begin{verbatim}
                docker compose logs -f alertas
                \end{verbatim}
                Respuesta: Logs indicando que la alerta fue procesada, e.g., \texttt{Received patient update: Test Patient}.
        \end{itemize}
\end{enumerate}

\section{Conclusión}
El sistema hospitalario basado en microservicios se implementó con éxito, cumpliendo los requisitos mínimos de tener al menos dos microservicios funcionales (Pacientes y Citas), con funcionalidad parcial para Alertas debido a problemas iniciales con la configuración de Kafka y Keycloak. Se resolvieron errores relacionados con el ID de clúster de Kafka y la conexión de Keycloak con PostgreSQL. Las pruebas confirman que el sistema es funcional y seguro, con JWT asegurando el acceso autorizado a los endpoints. Futuras mejoras incluyen la implementación de pruebas unitarias más exhaustivas, optimización para entornos de producción y resolución de problemas residuales con el microservicio Alertas.

\section{Referencias}
\begin{itemize}
    \item Repositorio GitHub: \url{https://github.com/Santos-Arellano/Software-System-Planning}
    \item Video de Demostración: \url{https://youtu.be/aZjE647RC1s}
\end{itemize}

\end{document}